% plan_v2.tex -- the plan for paper_v2: cutting the paper to about 8 pages of main text in ACM's
% two-column format. A standalone document (pdfLaTeX). The full version and its own plan
% (paper_plan.tex) stay in ../paper/ and are not edited for this one.
\documentclass[11pt,a4paper]{article}

\usepackage[T1]{fontenc}
\usepackage[utf8]{inputenc}
\usepackage{lmodern}
\usepackage[margin=2.2cm]{geometry}
\usepackage{microtype}
\DisableLigatures[-]{encoding = T1, family = tt*}   % -- stays two hyphens in commands
\usepackage{array}
\usepackage{booktabs}
\usepackage{tabularx}
\usepackage{enumitem}
\usepackage[hidelinks]{hyperref}

\DeclareRobustCommand{\code}[1]{\texttt{\detokenize{#1}}}
\newcolumntype{L}{>{\raggedright\arraybackslash}X}
\setlist{itemsep=2pt,topsep=3pt}
\setlength{\parskip}{4pt}
\setlength{\parindent}{0pt}
\setlength{\emergencystretch}{3em}

\title{Plan for \texttt{paper\_v2}\\[4pt]\large The 8-page version in ACM's two-column format}
\author{}
\date{7 November 2026}

\begin{document}
\maketitle

\section*{Status}

C1 to C8 are done (7 and 8 November). C1 set up \code{paper_v2/}; the author's answers to its
questions are decisions V1--V4 (Section~3), V5 was added after C2, V6 after C3, V7 after C4, V8
after C5, V9 after C6 and V10 after C7. C2 moved Tables~2 and~6 and the random draws of Table~3 to the appendix and made
Figures~2 and~4 one column wide (the numbers of that time; since C6, Figure~1 is in the appendix
as Figure~4 and the others are Figures~1--3): the main text went from about 11.5 to about 10.3
pages, with no argument cut. C3 cut the abstract and Sections~1--3 from 2,199 to 1,634 words, C4
cut the prose of Sections~4--6 from 3,477 to 2,610 words, C5 that of Section~7 from 1,731 to 1,097,
and C6 moved Figure~1 to a new Appendix~A and cut Sections~8--10 and the end statements from 1,489
to 1,056 words, moving details to the appendices and the settings list to Table~8. After C6, the main text,
from the abstract to the conclusion, ended about 7.96 pages in, within 8 pages, and the end
statements about 8.23 pages in. Figure and appendix numbers in this plan are those of their
time: C6's new Appendix~A moved the earlier Appendices~A--F to B--G, and C7 swapped the numbers of
Tables~6 and~7. C7 put the appendix in the paper's order: A the study at a glance, B the data and
the graph, C the changes from the base papers, D the search in detail, E the random draws and the
settings, F the full results, G reproducibility. C8 closed the last 0.23 page: with about 150
words cut or moved to the appendix (the abstract kept as V6 says), the code statement shortened as
V9 allows, and Figures~1--3 set a little lower, the paper from the abstract to the acknowledgments
now fills exactly 8 pages, and the references begin at the top of page~9. A review of the whole
version on 8 November changed three sentences (one each in Related work, Results and Appendix~E,
for clarity) and a few source comments that named old appendix letters; and, at the author's
choice, ten data shares that printed as whole percentages now have one decimal, as the rounding
rule asks (\code{make_numbers.py}, committed as \code{0bf5bc1}, which \code{numbers.tex}
now records). After the HAN repository's commits \code{a9141af} and \code{2ea02f6}, which changed
which results the recommendation tools run, Appendix~G says so, the ``More paths'' paragraph of
Appendix~F names HAN, as the full version's sentence does (for metapath2vec, neither size wins),
and \code{numbers.tex} records \code{2ea02f6}; no other value changed. The work
on this version is complete until the venue is chosen.

On 25 November, after the author's review of the Vietnamese version (\code{paper_v2_vn/}, part
L10), parts B1--B5 were planned to bring that review into this version
(Section~\ref{sec:backport}). The author decided D1--D6 there the same day, and B1--B5 are
done: the back-port is complete. A final review of both versions followed the same day
(Section~\ref{sec:final}).

\section{The target and why}

About \textbf{8 pages of main text} (from the abstract to the conclusion) in ACM's two-column
format (\code{acmart}, \code{sigconf}), plus up to 2 pages of references; the appendix becomes
supplementary material. Eight pages is the strictest limit among the likely venues, so one
version fits all of them:
\begin{itemize}
  \item RecSys, reproducibility papers (2026 rules): 8 pages plus up to 2 for references
        (\url{https://recsys.acm.org/recsys26/camera-ready/});
  \item SIGIR, reproducibility track (2023 rules, the latest posted): 9 pages plus references
        (\url{https://sigir.org/sigir2023/?p=1167});
  \item CIKM, full papers (2026 rules): 10 pages including any appendix, plus 2 for references
        (\url{https://cikm2026.diag.uniroma1.it/full-research-papers/}).
\end{itemize}
HAN's paper has about 9 pages of main text. If the venue chosen allows more, a page is easily
added back; the venue's own rules, once known, override this target.

\section{Rules}

\begin{itemize}
  \item \code{paper/} is the full version (thesis, arXiv, a journal such as TORS) and is not edited
        for this one, except for the corrections of V5. \code{paper_v2/} has its own \code{sections/}, \code{tables/},
        \code{figures/} and \code{numbers.tex}; the three \code{.bib} files stay shared at the
        project root.
  \item Nothing is lost from the work: what leaves the main text moves to the appendix
        (supplementary material) or stays in the full version. No number, limitation or caveat is
        dropped from the paper as a whole.
  \item No new claims, and the claims keep the qualifiers of the full version (``mostly'', seed
        noise, the warm-up limit, the small schema). Every number stays \verb|\val{key}| from
        \code{numbers.tex}, made by the LMSPS repository's script:
        \code{python paper/make_numbers.py --lmsps . --han ../han-improve --out <paper_v2>}.
  \item Each part ends with a clean compile (0 warnings) and a page count, and sends only the files
        it changed.
\end{itemize}

\section{Decisions}

\begin{enumerate}[label=V\arabic*]
  \item \textbf{Target (7 November):} about 8 pages of main text, up to 2 pages of references,
        and the appendix as supplementary material (Section~1).
  \item \textbf{The full version stays (7 November):} \code{paper/} keeps both its main files,
        \code{paper.tex} (one column) and \code{paper_acm.tex} (the full text in ACM's format,
        useful for arXiv), unchanged by this work.
  \item \textbf{Shared bibliographies (7 November):} the three \code{.bib} files stay at the
        project root for both versions, so a correction, such as the publisher details that
        ACM's style asks for (C8), shows in both reference lists.
  \item \textbf{File name (7 November):} the main file is \code{paper_v2/paper_v2.tex}.
  \item \textbf{Fixes go into both versions (7 November):} an error found during the cut is
        corrected in \code{paper/} as well, the one exception to V2. The first: Setup and Threats
        said every table shows popularity; it appears beside the models in the main comparison,
        the user-group table and HAN's full results only. C3 found two more. The abstract's last
        sentence now says that it holds for HAN, as the full version's claims do (in
        metapath2vec the manual sets beat the average random draw). And Related work and Threats
        said full ranking is reported for every result, while the warm-up check reports sampled
        Hit@10 only; they now say every path set. C4 found one more, in the ACM layout only:
        acmart ends a run-in heading with a full stop unless it ends with one already, and a stop
        after a capital does not count, so the headings LMSPS, HAN and GAT printed two stops;
        they now end with \verb|\@.| instead of a plain stop. C5 found three more, all in figure
        text. Figure~3's caption said that the main sets' three seeds ``differ by at most''
        their standard deviation, which understates how far they spread; it now says standard
        deviation. Its description for screen readers called it a bar chart, but it is a dot
        plot. And in this version only, Figure~2's description still said left and right after
        C2 had stacked its panels. C6 found one more, in the full version only, since this
        version's discussion no longer has the sentence: it said that the manual sets do as well
        as an average random draw, where Results says about as well. C8 changed two shared
        things that were not errors: the publisher, address and page fields of the shared
        bibliographies (V3), and the sides of the \code{PCP} and \code{PPP} labels, set in
        \code{make_numbers.py}, which also move in the full version's figure. The review after C8
        changed a third: ten data shares, such as the 75\% of the test edges on books and music,
        printed as whole percentages against the rounding rule and now have one decimal.
  \item \textbf{The front (7 November, the author's answers after C3):} this version's
        introduction has no roadmap paragraph, which the full version keeps. Related work leaves
        out the one-line accounts of how the cited methods work (HGRec, FHRec, AMERec, Zhu et
        al.'s path ranking, GEMS's predictor, RMS's reward), which stay in the full version, and
        its spare words go to Sections~4--8. The abstract keeps its
        length, 19 words over its target, until the venue's cap is known.
  \item \textbf{Making up the gap (8 November, the author's answer after C4):} Sections~1--6 end
        about 260 words over their budget, which would leave the paper about 0.3 page over 8.
        C5 and C6 therefore cut deeper than first planned: Results to about 1,600 words instead
        of 1,850 (for example, the user-group and product-group details to the appendix) and
        Discussion to about 430 instead of 450, with every number and caveat kept in the paper.
        If the paper is still over 8 pages after C6, C8 moves Figure~1, the overview of the
        study, to the appendix.
  \item \textbf{Figure~1 moves in C6 (8 November, the author's answer after C5):} a scratch
        build showed that the paper cannot reach 8 pages without V7's fallback unless C6 cuts
        about 300 words more than its targets. So C6, not C8, moves Figure~1, the overview of
        the study, to the appendix, and then cuts what the measured gap still needs, about 100
        words beyond its targets. The introduction's pointer to Figure~1 then names the
        appendix.
  \item \textbf{The last 0.23 page (8 November, the author's answer after C6):} the gap left after
        C6 is the length of the end statements, and Sections~8--10 are at about their targets. So
        C8 trims about 220 words from the sections that ended over theirs (the setup, Results,
        the data section, the introduction and the abstract), moving details to the appendix and
        keeping every number and caveat. Only if that is not enough does C8 shorten the code
        statement to a pointer to the reproducibility appendix; the AI statement stays as it is.
        Whether the chosen venue counts the end statements towards its limit is checked once the
        venue is known.
  \item \textbf{The appendix (8 November, the author's answers after C7):} it keeps the compact
        layout, its tables following its text in numbered order, unless the venue asks otherwise;
        page breaks that keep each table beside its section would cost 2--4 mostly empty pages.
        And it stays in the same PDF, after the references, until the venue is chosen; if the
        venue wants supplementary material as a separate file, it is split off then, keeping its
        letters so that the main text's references still hold.
\end{enumerate}

\section{The baseline and the budget}

Measured on 7 November in this folder's ACM layout. Before C2: 10,363 words from the abstract to
the conclusion, about 11.5 pages. After C2: 9,997 words, about 10.3 pages, plus the two end
statements and the acknowledgments (311 words, about 0.3 page), which count towards the limit at
most venues. The popularity fix (V5) added 10 words, so C3 started from 10,007. After C3: 9,443
words; the main text ends about 9.9 pages in, the end statements about 10.2. After C4: about
8,570 words; the main text ends about 9.0 pages in, the end statements about 9.4. After C5: about
7,940 words; the main text ends about 8.5 pages in, the end statements about 8.9. After C6, with
Figure~1 in the appendix: about 7,440 words; the main text ends about 7.96 pages in, the end
statements about 8.23. A word count includes
the text of the tables and figures that sit on a section's pages, so it is a guide, not an exact
share; after C4, Table~2 moved onto the first page of Results, so the C4 column gives each
section's words as if the floats had stayed (the count before C4 minus the prose cut). At about
970 words a page, 8 pages with the statements hold about 7,400 words of main text; measured after
C4, a page holds about 940. V7 lowers the targets of Results and Discussion so that the
total, with Sections~1--6 as they are, stays at about 7,380 words. Results ended 153 words over
its target. Moving Figure~1 to the appendix, V7's fallback, tried in a scratch build after C5,
brings both ends about 0.25 page earlier; with that, C6 has to cut about 100 words more than its
targets for the end statements to finish within 8 pages. V8 moves that step from C8 to C6. C6
brought Sections~8--10 and the end statements to about their targets, not 100 words below them:
cutting further would have removed arguments from the discussion. About 0.23 page, the length of
the end statements, is left to C8. After C8, the paper from the abstract to the acknowledgments
fills exactly 8 pages and the references begin at the top of page~9; the table does not track C8's
small cuts by section.

{\footnotesize
\begin{tabularx}{\textwidth}{@{}lrrrrrrL@{}}
  \toprule
  & \multicolumn{5}{c}{After} & & \\
  \cmidrule(lr){2-6}
  Section & C2 & C3 & C4 & C5 & C6 & Target & Mainly by \\
  \midrule
  Abstract (with CCS and keywords) & 248 & 219 & 219 & 219 & 219 & 200 & one sentence fewer \\
  1 Introduction & 864 & 670 & 670 & 670 & 673 & 650 & no findings repeated from the abstract; a
    shorter contribution list \\
  2 Related work & 719 & 495 & 495 & 495 & 495 & 550 & tighter paragraphs \\
  3 Preliminaries & 369 & 250 & 250 & 250 & 250 & 250 & shorter definitions \\
  4 Data & 1,184 & 1,184 & 916 & 916 & 916 & 850 & the twin-product check to the appendix; a
    shorter McAuley comparison \\
  5 Method & 1,797 & 1,797 & 1,448 & 1,448 & 1,306 & 1,350 & shorter vote and leakage paragraphs;
    Figure~1 to the appendix (V8) \\
  6 Experimental setup & 1,205 & 1,211 & 961 & 961 & 961 & 850 & the settings list to the appendix
    (Table~8 has it) \\
  7 Results & 2,387 & 2,387 & 2,387 & 1,753 & 1,753 & 1,600 & recommendation popularity, the
    $\delta$ check, user and product groups to the appendix; shorter cost quotations (V7) \\
  8 Discussion & 704 & 704 & 704 & 704 & 442 & 430 & ``Two metrics'' merged; tighter paragraphs
    (V7) \\
  9 Threats to validity & 361 & 366 & 366 & 366 & 294 & 260 & a tighter list \\
  10 Conclusion & 159 & 160 & 160 & 160 & 137 & 130 & \\
  \midrule
  Total & 9,997 & 9,443 & 8,576 & 7,942 & 7,446 & 7,120 & \\
  End statements and acknowledgments & 311 & 311 & 311 & 311 & 235 & 220 & C6 \\
  \bottomrule
\end{tabularx}\par}

\section{The parts}

{\small
\begin{tabularx}{\textwidth}{@{}>{\raggedright\arraybackslash}p{0.2\textwidth}L>{\raggedright\arraybackslash}p{0.15\textwidth}@{}}
  \toprule
  Part & What it does & Status \\
  \midrule
  C1 Set up & \code{paper_v2/}: its own copy of the ACM version, compiling on its own; this plan
    & done, 7 November \\
  C2 Tables and figures & 4 figures and 4 tables in the main text; Figures~2 and~4 one column
    wide; Tables~2 and~6 and the random draws of Table~3 to the appendix (a new Appendix~A);
    about 1.25 pages saved & done, 7 November \\
  C3 Front & abstract, Sections 1--3 to the budget: 2,199 to 1,634 words; the roadmap
    paragraph dropped; two claims corrected in both versions (V5) & done, 7 November \\
  C4 Data, method, setup & Sections 4--6 to the budget: prose 3,477 to 2,610 words, 275 over
    the budget; cluster sizes, the McAuley figures, the twin check, the title vocabulary and the
    ancestor links to Appendix~A, the runs' agreement and the vote's null to a new Appendix~C,
    the settings list to Table~8; no number dropped; one ACM heading fix (V5) & done,
    7 November \\
  C5 Results & Section 7 to the budget, about 1,600 words (V7): prose 1,731 to 1,097 words,
    1,753 with its floats; ``What the search finds'' merged into ``What HAN does with the paths'';
    the paths' ranks and the $\delta$ check to Appendix~C; the user groups, more paths, the path
    weights and preferences, and HAN against metapath2vec to Appendix~E; no number dropped; three
    fixes in figure text (V5) & done, 8 November \\
  C6 Discussion and end & Figure~1 to a new Appendix~A (V8); Sections 8--10 and the two end
    statements from 1,489 to 1,056 words: ``Two metrics'' merged into the opening and then, with
    its details, into Appendix~F; the threats that the setup already states shortened; the list of
    recommendation tools to Appendix~G; the paragraph on the user side kept; no number dropped
    & done, 8 November \\
  C7 Supplementary material & everything moved in C2--C6 found in the appendix; the random draws
    from the data appendix to the settings (now ``The random draws and the settings''); the 22
    references to an appendix checked against their targets; three missing pointers added (Table~10
    and the two metrics in Results, Tables~9 and~10 in the setup); the compact layout kept, since
    page breaks that keep each table beside its section cost 2--4 mostly empty pages; whether the
    appendix becomes a separate PDF waits for the venue (V10) & done, 8 November \\
  C8 Final pass & 8 pages: about 150 words cut or moved to Appendices~E and~G (the setup, Results,
    the introduction, related work, the discussion; the abstract kept, V6), the code statement
    shortened to a pointer to Appendix~G (V9), Figures~1--3 set 0.6--1.6\,cm lower; the main text
    with the end statements and acknowledgments ends on page~8; every number still from
    \code{numbers.tex} and every typed number a design constant; the claims of every changed
    sentence checked against the full version; 0 LaTeX warnings, and BibTeX's 40 notices down to
    9 by publisher, address and page fields in the shared bibliographies (V3), the 9 left having
    no value to give (ICLR and PMLR have no publisher address, papers without page numbers, one
    page count that only the published TOIS copy would show); the labels of \code{PCP} and
    \code{PPP} moved apart in \code{make_numbers.py}, which the author committed as
    \code{1b90835}: run from that commit, the script reproduces every figure file, and
    \code{numbers.tex} changes only in the commit it records; the package & done, 8 November \\
  \bottomrule
\end{tabularx}\par}

\section{Back-port of the Vietnamese review (B1--B5)}\label{sec:backport}

\textbf{Why.} From 20 to 22 November the author reviewed the Vietnamese version by hand
(\code{paper_v2_vn/}, part L10, decisions N17--N32 in its plan). This version has not changed since
10 November, so it now takes over that review, in five parts of one request each. The source is
the final \code{paper_v2_vn}: the author's own build of 22 November, which matches the assistant's
copy file for file. The Vietnamese text before L10 is a sentence-by-sentence translation of this
version, so comparing the two Vietnamese texts shows which English sentences change. L10 changed
30 files, with about 900 lines in and 800 out.

\textbf{What is ported and what is not.}
\begin{itemize}
  \item \emph{Ported:} every change of fact, claim, framing, structure, number, citation or
        cross-reference. The main ones: search compared with manual design, with the random sets as
        the control (N29); a verdict that depends on the context (N30); the abstract and the
        introduction as the author reviewed them (N18, N21, N22); the related-work changes (N27);
        no Discussion or Threats section (N26); no CCS concepts (N19); no appendix but a
        supplementary document, with the parts of the appendix that moved into the text (N31).
  \item \emph{Not ported:} what concerns Vietnamese only. That covers the terms for meta-path
        (N17), random sets, embeddings, path weights, path instances and input features
        (N22--N25), ancestor categories (N32) and a repository, and the comma before ``and'' (N20).
  \item \emph{Decided by the author:} D1--D6 below.
\end{itemize}

\textbf{How.}
\begin{enumerate}
  \item For each file, compare the Vietnamese before L10 with the final Vietnamese. Each changed
        Vietnamese sentence points to the English sentence it was translated from.
  \item A change of substance is written into the English, in the English paper's style and with
        its wording wherever the meaning did not change. A Vietnamese-only change is left out.
  \item Every number stays a \verb|\val| from \code{numbers.tex}, which needs no rebuild: it has the
        same 1,262 keys as the Vietnamese one, including the 518 that the final Vietnamese uses.
        Each citation keeps the scope checked in L10.
  \item Each part ends with a pdfLaTeX build with 0 warnings, a check that the English and the
        Vietnamese say the same things, and a package of the changed files and the PDF only.
\end{enumerate}

\textbf{Decisions (the author, 25 November).}
\begin{enumerate}[label=D\arabic*]
  \item \emph{Title:} the descriptive title, as in the Vietnamese (N17): ``Cost and Effectiveness
        of Meta-Path Search for HAN-Based Product Recommendation''. The subtitle stays.
  \item \emph{Semicolons and close commas:} the author's rules for the Vietnamese hold for the
        English too. Its 150 semicolons (text, tables, figures and PDF keywords) go, as in N20, and
        no new sentence puts two commas close together.
  \item \emph{AI statement:} the English follows the author's Vietnamese wording (N28). It no
        longer says that the assistant helped draft and edit the text, and it says that the
        assistant helped write the script.
  \item \emph{CCS concepts:} removed, as in N19. The keywords follow the Vietnamese, with ``manual
        meta-path design'' instead of ``random baselines''.
  \item \emph{Appendix:} a supplementary document, \code{supplement_v2.tex}, as V1 and N31 say,
        which the paper does not refer to.
  \item \emph{Length:} not a concern in B1--B5, which translate. The author looks at the length
        afterwards, and B4 only reports where the main text ends.
\end{enumerate}

{\small
\begin{tabularx}{\textwidth}{@{}>{\raggedright\arraybackslash}p{0.14\textwidth}L>{\raggedright\arraybackslash}p{0.28\textwidth}@{}}
  \toprule
  Part & What it ports & Files \\
  \midrule
  B1 Front & the title (D1), the CCS concepts and the keywords (D4); the abstract (N18, N29, N30);
    Section~1 (N21, the author's introduction of N22, N29, N30); Section~2 (N27, and the new
    paragraph on how a search method is judged, N29)
    & \code{paper_v2.tex}, \code{sections/abstract.tex}, \code{intro.tex}, \code{related.tex} \\
  B2 Sections 3--6 & Section~3; Section~4 (the input features, the McAuley data page as a footnote
    from the appendix, no appendix pointers); Section~5 (the author's final text: no opening
    pointers, the runs' agreement in the text, the one-standard-error rule and the warm-up as
    reworded); Section~6 (three kinds of sets with the random sets as the control, the rank, why
    Recall@10 and Precision@10 are not reported, seeds and statistics, the key settings and
    stopping epochs, the software); Table~3 with the ten random sets side by side and HAN's
    validation pick marked; Table~1 and Figure~1. The largest part.
    & \code{sections/prelim.tex}, \code{data.tex}, \code{method.tex}, \code{setup.tex},
    \code{tables/t1_data.tex}, \code{t3_sets.tex}, \code{figures/f2_ties.tex} \\
  B3 Results & Section~7: 7.1 (\code{a2} and \code{c} with Figure~2, the range over user groups);
    7.2 renamed for the random control (the cheapest choice, HAN against metapath2vec); 7.3
    renamed for cost against effectiveness (the author's paragraph, the search's GPU hours as the
    yardstick); 7.4 with why the searched sets lose (from Section~8, marked untested) and the
    correlation within the random sets; 7.5 with why GAT leads (from Section~8) and its lead by
    user group; 7.6 as reworded; Tables~4 and~7, Figures~2 and~3
    & \code{sections/results.tex}, \code{tables/t4_main.tex}, \code{t7_cost.tex},
    \code{figures/f3_sets.tex}, \code{f4_weights.tex} \\
  B4 Conclusion and end & Sections~8--10 become one Conclusion (N26, N30): the verdict that depends
    on the context, the lesson that a search should be judged by the recommender it serves, the
    limitations with the two minor ones, future work. The code statement as in the Vietnamese, the
    AI statement (D3). The paper without its appendix (D5), the last page checked, every
    cross-reference resolved, the length reported (D6)
    & \code{sections/discussion.tex}, \code{statements.tex}, \code{paper_v2.tex} \\
  B5 Supplement and final pass & \code{supplement_v2.tex} from the appendix: sections A--H with the
    new ``Other limitations'', tables S1--S7 and figure S1, the paper's numbers through the
    \code{xr} package. The appendix's L10 changes and those of Tables~2, 3b, 5, 6, 8, 9 and
    Figure~4. A \code{latexmkrc} for pdfLaTeX with the \code{xr} rule and the check for the
    \code{.bib} files. Both documents with 0 warnings, a last check of every number, citation and
    claim against the Vietnamese, this plan, and the full folder with the three \code{.bib} files
    and both PDFs
    & \code{supplement_v2.tex} and \code{latexmkrc} (new), \code{sections/appendix.tex},
    \code{tables/t2_graph.tex}, \code{t3b_random.tex}, \code{t5_groups.tex}, \code{t6_changes.tex},
    \code{t8_settings.tex}, \code{t9_full.tex}, \code{figures/f1_pipeline.tex}, \code{plan_v2.tex} \\
  \bottomrule
\end{tabularx}\par}

\textbf{B1, done 25 November.} The title is now the descriptive one, broken after ``Search'' so
that ``HAN-Based Product Recommendation'' stays on one line. The CCS concepts are a comment that
says how to restore them, the keywords and PDF keywords end with ``manual meta-path design'', and
the PDF keywords are separated by commas. The abstract, Section~1 and Section~2 say what the
reviewed Vietnamese says. Their semicolons are gone, and the new sentences put no two commas close
together. Two old sentences with a short aside between commas now use parentheses (``paths (or
richer structures)'', ``ENAS (a leading search method)''). The English keeps ``one-shot'' for
LMSPS and DiffMG's kind of search, which the Vietnamese dropped only because its own words read as
a search run once, and keeps Yu et al.'s quotation, which the Vietnamese paraphrased to avoid
quoting English. The sentence on time-ignoring splits left Section~2 for the conclusion (B4). Its
two citations stay in Section~9 until then, so the 25 references are unchanged. \code{LMSPS} is no
longer hyphenated (a \verb|\hyphenation| entry). The build has 0 warnings and 14 pages, and the
main text still ends 8.02 pages in.

\textbf{B2, done 25 November.} Sections~3--6 say what the reviewed Vietnamese says. There are no
appendix pointers, and the McAuley data page is now the footnote that gives the quotation's source.
The input features are described as the author did (the product's group and LSA title vector, the
user's hashed binary vector). Section~5 follows the author's rewording: the leakage controls, the
one-standard-error window, the runs' agreement in the text, the last sentence on the six-path
choice and the reason for the warm-up. Section~6 compares the searched sets with the manual ones,
with the random draws as the control. It numbers the manual principles, spells out the rank, says
why Recall@10 and Precision@10 are not reported, and gives the key settings, the stopping epochs
and the software. As in the Vietnamese, its manual item no longer says ``with AI assistance'',
which the introduction, the sets table and the AI statement still say. The sets table lists all 24
sets, with the random draws side by side and HAN's validation pick marked. Two English-only choices
were made. ``Sampled'' is used for the negatives and LMSPS's paths, so that ``draws'' means the
random sets, as the Vietnamese does with its own words. GAT is a ``baseline'' in Section~5, where
the Vietnamese still says control. Section~4's pointer for the split's limitation stays on
Section~9 until B4 removes that section. These files have no semicolons, and the new sentences put
no two commas close together. The build has 0 warnings and 15 pages, and the main text ends 8.20
pages in.

\textbf{B3, done 25 November.} Section~7 says what the reviewed Vietnamese says. 7.2 is now ``The
random control''. It places the searched and the manual sets among all 24 with the random draws as
the control, and it gives the cheapest choice (from Section~8) and HAN against metapath2vec (a
Spearman correlation from the appendix). 7.3 is now ``Cost against effectiveness'', the author's
paragraph. The hand design's time was not recorded, so the search's GPU hours are the yardstick:
hand design is cheaper while it takes less, and the search has not bought HAN better results. 7.4
ends with why the searched sets lose (from Section~8, with the caveat that the check only locates
the loss) and gives the correlation within the random draws. 7.5 explains why GAT leads (from
Section~8), with its lead by user group. 7.1 gives the range over user groups, and no appendix
pointer is left in Section~7. Tables~3 and~4 and Figures~2 and~3 have no semicolons now. Section~8
still holds the moved paragraphs until B4 merges Sections~8--10 into the conclusion, so the main
text ends 8.55 pages in for now. The build has 0 warnings and 15 pages.

\textbf{B4, done 25 November.} Sections~8--10 are now one Conclusion (Section~8), as in the reviewed
Vietnamese. It gives the verdict that depends on the context. Manual design is the better choice
when an expert can design the paths in less time than the search's roughly 20 GPU hours. Without an
expert, search gives sets close to manual design. The random control does at least as well as both,
and any advantage of search more likely lies in larger schemas. The conclusion then gives the lesson
that a search should be judged by its recommender, a Limitations paragraph with the two minor
limitations of the old Threats list, and the future work as a numbered list. The other Threats
items are covered by Sections~6 and~7.4 or go to the supplementary document. The code statement
gives the commits, the data rebuild, the run records and \code{make_numbers.py}. The AI statement
follows the author's Vietnamese wording (D3). The paper has no appendix, Section~4's pointer for
the split now goes to the conclusion, and no reference to the Discussion, the Threats or the
appendix is left. Column balancing stays off, as in the Vietnamese, since acmart's
\verb|\balance| would warn. Length (D6): the main text from the abstract to the conclusion ends
about 7.9 pages in, within V1's 8 pages. The end statements and acknowledgments end 8.18 pages in,
and the references fill page~9. The whole main paper has no semicolons, its 25 references are
unchanged, and the build has 0 warnings and 9 pages.

\textbf{B5, done 25 November.} The appendix is now the supplementary document,
\code{supplement_v2.tex}, which has the paper's preamble. It has sections A--H (H, other
limitations, is new), tables S1--S7 and figure S1, and it prints the paper's numbers through the
\code{xr} package (labels prefixed \code{P-}, printed with \verb|\ref*|). It holds the reviewed
Vietnamese supplement's changes: the overview figure with the random draws as the control and GAT
and popularity as baselines, the clearer sentences of the reproducibility section, and no
semicolons, as in the rest of the English. A \code{latexmkrc} for pdfLaTeX builds both documents,
with the \code{xr} rule and the check for the \code{.bib} files. From an empty folder,
\code{latexmk} builds both, and \code{latexmk supplement_v2.tex} builds the paper first.
\code{numbers.tex} now records the LMSPS commit \code{5e32eb6}, as the Vietnamese does. Run again
from the two repositories, \code{make_numbers.py} changes nothing else in it or in the figure data.
The HAN commit stays \code{2ea02f6}, as the two later commits change only the recommendation API.
A last check against the Vietnamese found that every file of the paper and the supplement uses the
same numbers, citations and cross-references, with the same headings and number of paragraphs.
The paper has 9 pages, with the main text ending about 7.9 pages in, and the supplement has 6. Both
build with 0 warnings and no overfull or underfull box. The back-port is complete.

\section{Final review (25 November)}\label{sec:final}

After B5 the author answered four questions, and a last reading of both versions, papers and
supplements, followed. The answers:
\begin{enumerate}
  \item In the Vietnamese Section~5, GAT is now called a baseline, with the word used for it
        everywhere else. The Vietnamese word for the control now always means the random draws,
        except in the author's AI statement, which keeps its wording. Nothing changes in the
        English.
  \item The Vietnamese supplement now calls a repository a code base, as the code statement
        does, instead of the glossary's earlier word. The English supplement keeps
        ``repository''.
  \item Short asides between two commas are gone from both versions, in old text as well as
        new. Each became a parenthesis, a sentence of its own or a sentence without the aside,
        with the meaning unchanged. Lists keep their commas, and the AI statement is untouched.
        In the English: the introduction (\code{PUP}), Section~3 (the node types, the positive
        reviews), Section~4 (the record, \code{PUP}, the share of positive lines, \code{PCCCP},
        the 5-core, the split), Section~5 (the score, metapath2vec), Section~6 (the searched
        sets), Section~7.4 (\code{PUPUP}, the near-empty paths), Figure~3's caption, Table~S4,
        three figure descriptions and three sentences of the supplement. Related work now says
        that HGRec, FHRec and AMERec fuse their embeddings ``as HAN does'', since ``Like HAN,
        HGRec, FHRec and AMERec'' read like a list of four.
  \item The supplements no longer print the branch names, since the commit hashes identify the
        code. The helper \verb|\codeval| that set them is gone from \code{supplement_v2.tex}.
\end{enumerate}
The reading found nothing else to change in the English. Every number is still a \verb|\val|,
the two versions use the same numbers, citations and cross-references in every file, and the
visible text has no semicolon. Both documents build with 0 warnings and no overfull or underfull
box: the paper has 9 pages, with the main text ending about 7.9 pages in, and the supplement 6.

\end{document}
